% =====================================================================
% SECCIÓN: METODOLOGÍA
% Simulador de Entrevistas de Trabajo — ISC
% Instituto Tecnológico Superior de San Pedro de las Colonias
% =====================================================================

\section{Metodología}

% ---------------------------------------------------------------------
\subsection{Análisis de requerimientos}
% Cronograma: 31 ago – 13 sep

Para el desarrollo del simulador se identificaron los requerimientos
funcionales y no funcionales del sistema a partir de los objetivos
planteados en el protocolo de investigación.

\subsubsection{Requerimientos funcionales}

Los requerimientos funcionales describen las acciones que el sistema
debe ser capaz de realizar:

\begin{enumerate}
    \item El sistema debe permitir al usuario crear una cuenta e
    iniciar sesión con correo electrónico y contraseña.

    \item El sistema debe permitir al usuario configurar su sesión
    de práctica eligiendo el área técnica de enfoque, el tipo de
    preguntas y la cantidad a responder.

    \item El sistema debe seleccionar y presentar preguntas de
    entrevista de forma ordenada por fases, respetando el flujo
    natural de una entrevista real: apertura, preguntas conductuales,
    preguntas técnicas y cierre.

    \item El sistema debe permitir al usuario responder las preguntas
    de dos formas: escribiendo su respuesta o respondiendo por voz.

    \item El sistema debe enviar cada respuesta a un modelo de
    inteligencia artificial y mostrar al usuario una evaluación
    con puntuación, fortalezas, áreas de mejora, un consejo
    específico y un ejemplo de respuesta ideal.

    \item El sistema debe guardar el historial de sesiones y
    respuestas de cada usuario para su consulta posterior.

    \item El sistema debe mostrar estadísticas de progreso al usuario:
    sesiones realizadas, preguntas respondidas y puntuación promedio.

    \item Al finalizar cada sesión, el sistema debe mostrar una
    pantalla de resultados con la puntuación general, la mejor y
    peor respuesta de la sesión, y un consejo global generado por
    la inteligencia artificial.
\end{enumerate}

\subsubsection{Requerimientos no funcionales}

Los requerimientos no funcionales describen las condiciones bajo las
cuales el sistema debe operar:

\begin{enumerate}
    \item El simulador debe funcionar desde un navegador web sin
    necesidad de instalar ningún programa adicional.

    \item La interfaz debe estar completamente en español y orientada
    al contexto del mercado laboral mexicano en el sector tecnológico.

    \item El sistema debe ser accesible desde los equipos del
    laboratorio de sistemas del ITSPC con conexión a internet.

    \item El tiempo de respuesta de la evaluación por inteligencia
    artificial no debe superar los cinco segundos en condiciones
    normales de conexión.

    \item El sistema debe proteger los datos de cada usuario de forma
    que ningún participante pueda acceder a la información de otro.

    \item En caso de fallo del servicio de inteligencia artificial,
    el sistema debe continuar funcionando con una evaluación básica
    para no interrumpir la sesión del usuario.

    \item El sistema debe evitar mostrar preguntas que el usuario
    ya haya respondido en sesiones recientes, garantizando variedad
    en la práctica.
\end{enumerate}

\subsubsection{Recursos tecnológicos}

Para cumplir con los requerimientos anteriores se seleccionaron las
siguientes herramientas y tecnologías:

\begin{itemize}
    \item \textbf{TypeScript:} Lenguaje de programación principal
    del proyecto, basado en JavaScript con tipado estático.

    \item \textbf{React:} Biblioteca utilizada para construir la
    interfaz de usuario del simulador.

    \item \textbf{Tailwind CSS:} Se empleó para los estilos visuales
    y el diseño responsivo de la aplicación.

    \item \textbf{Supabase:} Base de datos en la nube (PostgreSQL)
    que almacenó los usuarios, las sesiones de práctica y las
    respuestas. También gestionó el inicio de sesión y el registro
    de usuarios mediante autenticación con correo y contraseña.

    \item \textbf{Google Gemini API:} Modelo de lenguaje de gran
    escala utilizado para evaluar las respuestas del candidato y
    generar retroalimentación personalizada en tiempo real.

    \item \textbf{Web Speech API:} Tecnología nativa del navegador
    que permitió capturar la voz del usuario y convertirla en texto
    para el modo de entrevista por voz. No requirió instalación ni
    costo adicional.

    \item \textbf{Bolt (bolt.new):} Plataforma utilizada para generar
    el prototipo funcional inicial del simulador a partir de una
    descripción en lenguaje natural.

    \item \textbf{Kiro IDE:} Entorno de desarrollo con agente de
    inteligencia artificial donde se desarrolló, refinó y depuró
    el código del proyecto.
\end{itemize}

Para la construcción del banco de preguntas se revisaron plataformas
especializadas en entrevistas de trabajo del sector tecnológico,
entre ellas: DataCamp, FinalRoundAI, ResuFit, InterviewPrep.org,
TestGorilla, igotanoffer y Primly.io. El resultado fue un banco de
70 preguntas curadas en español, organizadas por tipo y nivel de
dificultad, orientadas al perfil de egreso de Ingeniería en Sistemas
Computacionales.

% ---------------------------------------------------------------------
\subsection{Diseño funcional}
% Cronograma: 7 sep – 20 sep

El diseño del simulador abarcó dos aspectos principales: la estructura
de la base de datos y el diseño de la interfaz de usuario.

\subsubsection{Diseño de la base de datos}

Con base en los requerimientos identificados, se diseñó el esquema
de la base de datos en Supabase utilizando PostgreSQL. Se definieron
cuatro tablas principales:

\begin{itemize}
    \item \textbf{profiles:} En esta tabla se almacenó la información
    personal de cada usuario registrado, como el nombre completo,
    el área profesional de interés y el nivel de experiencia.

    \item \textbf{interview\_questions:} Aquí se almacenaron las
    70 preguntas del banco, cada una clasificada por fase (apertura,
    conductual, técnica o cierre), área técnica (general, web,
    backend o bases de datos), nivel de dificultad y criterios de
    respuesta ideal utilizados por la inteligencia artificial.

    \item \textbf{interview\_sessions:} En esta tabla se registró
    cada sesión de práctica iniciada por un usuario, guardando el
    modo de interacción elegido, el área técnica, el tipo de sesión
    y la cantidad de preguntas respondidas.

    \item \textbf{session\_responses:} Aquí se guardaron las
    respuestas del usuario por cada pregunta respondida, junto con
    la puntuación, la retroalimentación y el consejo de mejora
    generados por la inteligencia artificial.
\end{itemize}

\subsubsection{Diseño de la interfaz}

El diseño de la interfaz se realizó inicialmente con Bolt, que generó
las pantallas principales del simulador. Posteriormente se refinó en
Kiro IDE, ajustando la distribución de elementos, la paleta de
colores y la experiencia de usuario. El simulador se compone de
ocho pantallas principales:

\paragraph{Inicio de sesión y registro.}
Es la primera pantalla que ve el usuario al entrar al simulador.
Desde aquí puede crear una cuenta nueva ingresando su nombre, correo
y contraseña, o iniciar sesión si ya tiene una cuenta registrada.
Incluye un indicador visual de fortaleza de contraseña durante
el registro.

\paragraph{Dashboard.}
Es la pantalla principal del simulador. Muestra el nombre del usuario,
un resumen de su progreso (sesiones realizadas, preguntas respondidas
y puntuación promedio), una gráfica de barras con la evolución de
su puntuación por sesión y los últimos consejos recibidos de la
inteligencia artificial. Desde aquí el usuario elige si quiere
practicar con el modo de entrevista por voz o por texto.

\paragraph{Configuración de sesión.}
Antes de iniciar una entrevista, el usuario personaliza su sesión
en esta pantalla. Elige el área técnica en la que quiere enfocarse
(General ISC, Desarrollo Web, Backend o Bases de Datos), el tipo de
preguntas que desea practicar (técnicas, conductuales o mixtas) y
cuántas preguntas quiere responder (5, 10 o 15). Al final se muestra
un resumen de lo seleccionado antes de comenzar.

\paragraph{Entrevista por texto.}
En esta pantalla aparece una pregunta a la vez con su categoría,
nivel de dificultad y fase. El usuario escribe su respuesta en el
área de texto disponible y la envía para ser evaluada. Se muestra
un contador de palabras en tiempo real y una barra de progreso que
indica el avance dentro de la sesión.

\paragraph{Entrevista por voz.}
Funciona igual que la entrevista por texto, pero en lugar de
escribir, el usuario responde hablando. Al presionar el botón del
micrófono comienza la grabación con un indicador visual animado y
la transcripción aparece en pantalla en tiempo real. Cuando el
usuario termina, la respuesta transcrita se envía automáticamente
para su evaluación.

\paragraph{Retroalimentación.}
Después de evaluar cada respuesta, el simulador muestra los resultados
en esta pantalla. El usuario puede ver la puntuación obtenida del 1
al 10, las fortalezas de su respuesta, las áreas de mejora, un
consejo puntual y un ejemplo de cómo se vería una respuesta ideal.
Al terminar puede continuar con la siguiente pregunta o finalizar
la sesión.

\paragraph{Resultados de sesión.}
Al finalizar la sesión completa, se muestra una pantalla de resumen
con la puntuación promedio obtenida, la pregunta donde el usuario
respondió mejor, la pregunta donde tuvo más dificultad, el detalle
de puntuación por cada pregunta respondida y un consejo general
generado por la inteligencia artificial con base en el desempeño
de toda la sesión.

\paragraph{Perfil.}
Aquí el usuario puede actualizar su información personal como el
nombre, el área profesional y su nivel de experiencia. También puede
consultar el historial completo de sus respuestas anteriores con la
fecha, la puntuación y la retroalimentación recibida en cada sesión.

% ---------------------------------------------------------------------
\subsection{Desarrollo del simulador web}
% Cronograma: 14 sep – 4 oct

El desarrollo se llevó a cabo en Kiro IDE sobre el código base
generado por Bolt. Se implementaron los siguientes módulos:

\subsubsection{Autenticación y gestión de usuarios}

Se implementó el registro e inicio de sesión con correo y contraseña
mediante Supabase. Al momento del registro se creó automáticamente
un perfil en la base de datos con los datos del usuario. La sesión
se mantuvo activa entre visitas mediante tokens almacenados en el
navegador. Se agregó validación de fortaleza de contraseña en tiempo
real durante el registro.

\subsubsection{Configuración de sesión}

Se desarrolló la pantalla de configuración que apareció antes de
iniciar cada entrevista. El usuario seleccionó el área técnica de
su interés, el tipo de preguntas y la cantidad. Estas preferencias
se guardaron en la base de datos al iniciar la sesión y se utilizaron
para filtrar el banco de preguntas.

\subsubsection{Selección de preguntas por fases}

Se implementó un algoritmo que distribuyó las preguntas respetando
el orden natural de una entrevista real: una pregunta de apertura
al inicio, preguntas conductuales o técnicas en el cuerpo de la
sesión según el tipo elegido, y una pregunta de cierre al final.
Dentro de cada fase la selección fue aleatoria, garantizando variedad
entre sesiones. Adicionalmente, el algoritmo verificó el historial
reciente del usuario para evitar repetir preguntas ya respondidas
en las últimas sesiones, priorizando siempre preguntas nuevas.

\subsubsection{Módulo de entrevista por texto}

Se desarrolló la pantalla de entrevista por texto que presentó cada
pregunta con su categoría, nivel de dificultad y fase. El usuario
escribió su respuesta y la envió para su evaluación. Se mostró un
contador de palabras en tiempo real y una barra de progreso con el
avance de la sesión. Se implementó una confirmación antes de
abandonar la sesión para evitar pérdida accidental del progreso.

\subsubsection{Módulo de entrevista por voz}

Se implementó la entrevista por voz utilizando la Web Speech API
del navegador, configurada en español. Al presionar el botón del
micrófono se inició la grabación con un indicador visual de ondas
concéntricas animadas y un contador de tiempo. La transcripción
apareció en pantalla en tiempo real. Al detener la grabación, el
texto transcrito se envió automáticamente a evaluación. Se agregó
un aviso visible al inicio de cada pregunta recomendando al usuario
estar en un lugar silencioso para mejorar la precisión de la
transcripción.

% ---------------------------------------------------------------------
\subsection{Conexión con IA y sistema de voz}
% Cronograma: 21 sep – 4 oct

\subsubsection{Integración con Google Gemini API}

Se desarrolló el módulo de evaluación que integró Google Gemini
(modelo \textit{gemini-3.5-flash}) como motor de retroalimentación.
Para cada respuesta se construyó un prompt contextualizado que
incluyó los siguientes elementos:

\begin{itemize}
    \item La pregunta y los criterios de una respuesta ideal.
    \item La fase y categoría de la pregunta (apertura, conductual,
    técnica o cierre).
    \item El área técnica y el nivel de experiencia del candidato.
    \item El modo de respuesta (escrita o por voz).
    \item El tiempo que tardó el usuario en responder.
    \item El historial de preguntas y puntuaciones previas dentro
    de la misma sesión, para que la inteligencia artificial pudiera
    detectar patrones de mejora o declive.
    \item Detección automática de muletillas en respuestas por voz
    (palabras como ``este'', ``o sea'', ``pues'', ``um'').
\end{itemize}

Google Gemini devolvió una evaluación estructurada en formato JSON
con los siguientes campos: puntuación del 1 al 10, fortalezas de
la respuesta, áreas de mejora, un consejo específico para esa
pregunta y un ejemplo de cómo comenzaría una respuesta ideal. Esta
información se almacenó en la base de datos y se presentó al usuario
al finalizar cada pregunta.

Se implementaron medidas de seguridad sobre el módulo de IA:
sanitización de las respuestas del usuario antes de enviarlas al
modelo para prevenir inyecciones de prompt, límite de longitud de
3{,}000 caracteres por respuesta y un tiempo máximo de espera de 15
segundos para la llamada a la API. En caso de fallo del servicio,
el módulo activó automáticamente un sistema de respaldo con
evaluación básica local para no interrumpir la sesión del usuario.

\subsubsection{Pantalla de resultados de sesión}

Al finalizar cada sesión se implementó una pantalla de resultados
que mostró: la puntuación promedio de la sesión con indicador visual
de desempeño, la pregunta con mayor puntuación, la pregunta con
menor puntuación, el detalle numérico de cada pregunta respondida
y un consejo general generado por Google Gemini analizando el
patrón de respuestas de toda la sesión.

\subsubsection{Gráfica de progreso}

Se implementó una gráfica de barras en el dashboard que mostró la
evolución de la puntuación promedio del usuario a lo largo de sus
sesiones. La gráfica apareció automáticamente cuando el usuario
completó dos o más sesiones, con un indicador textual de tendencia
(mejorando, en declive o estable).

% ---------------------------------------------------------------------
\subsection{Pruebas y correcciones}
% Cronograma: 28 sep – 11 oct

Durante esta fase se realizaron pruebas funcionales del simulador
para identificar y corregir errores antes de la intervención con
los participantes. Las pruebas abarcaron los siguientes aspectos:

\subsubsection{Pruebas funcionales}

Se verificó el correcto funcionamiento de cada módulo del simulador:
registro e inicio de sesión, configuración de sesión, carga y
filtrado de preguntas, envío de respuestas, recepción y visualización
del feedback de la inteligencia artificial, pantalla de resultados
y consulta del historial en el perfil del usuario.

\subsubsection{Correcciones implementadas}

A partir de las pruebas se identificaron y corrigieron los siguientes
problemas:

\begin{itemize}
    \item Se corrigió un error en los enlaces del menú de navegación
    que permitía iniciar una entrevista sin pasar por la pantalla
    de configuración.

    \item Se agregó manejo de errores en el guardado del perfil
    del usuario, cambiando la operación de actualización por una
    operación de inserción o actualización (\textit{upsert}) para
    garantizar que los datos se guardaran correctamente en todos
    los casos.

    \item Se corrigió un error de tipado en el componente de perfil
    que impedía compilar el proyecto correctamente.

    \item Se implementó la confirmación antes de abandonar una
    entrevista en curso para evitar pérdida accidental de respuestas.

    \item Se verificó la conexión con la base de datos de Supabase
    y se validó que las políticas de seguridad por filas (RLS)
    funcionaran correctamente para cada usuario.
\end{itemize}

\subsubsection{Pruebas de usabilidad preliminares}

Se realizaron pruebas preliminares de uso por parte de los propios
desarrolladores para validar el flujo completo de una sesión de
práctica, desde el registro hasta la visualización de resultados,
tanto en el modo de entrevista por texto como en el modo por voz.
Se verificó el funcionamiento en los navegadores Google Chrome y
Microsoft Edge, que son los disponibles en los equipos del
laboratorio del ITSPC.

\end{document}
